\section{Versuchsbeschreibung}
Der Versuch zur Untersuchung der Solarzelle gliedert sich in drei aufeinander aufbauende Teile. Im ersten Teil wird die elektrische Kennlinie der Solarzelle aufgenommen, im zweiten Teil die spektrale Empfindlichkeit untersucht und im dritten Teil die Umwandlung und Speicherung der gewonnenen Energie anhand eines Hydro-Cars demonstriert. 

\subsection{Statische Solarzelle}
Im ersten Versuchsteils wird ein Solarzellenmodul untersucht. Dieses wird mit einer Halogenglühlampe, als künstliche Strahlungsquelle, beleuchtet und die Intenstät der einfallenden Strahlung mit einem regelbaren Transformator variiert. 
Die Solarzelle wird an einen Stromkreis mit einer Widerstandsdekade angeschlossen. Durch schrittweises Verändern des Widerstands kann der Arbeitspunkt gezielt verschoben werden. Zwei Multimeter erfassen dabei gleichzeitig Stromstärke und Spannung: Die Strommessung erfolgt in Reihe, während die Spannung parallel zur Solarzelle gemessen wird. Auf diese Weise werden für verschiedene Widerstandswerte Messpaare aufgenommen, aus denen sich die U-I-Kennlinie der Solarzelle ergibt. 
\begin{figure}[H]
    \centering
    \includegraphics[width=0.4\linewidth]{VersuchsaufbauMuster.png}
    \caption{Versuchsaufbau zur Aufnahme von U-I-Kennlinien einer Solarzelle.~\cite{leifi}} \label{fig:V1}
\end{figure}
\noindent Ergänzend werden die charakteristischen Größen Kurzschlussstrom $I_K$ und Leerlaufspannung $U_L$ bestimmt. Der Kurzschlussstrom wird durch ineinander Stecken der Kontakte und damit durch kurzschließen des Stromkreises gemessen, während für die Bestimmung der Leerlaufspannung der Stromkreis geöffnet wird, sodass kein Strom mehr fließt. 
Diese beiden Größen definieren die Randpunkte der Kennlinie und sind wesentlich für die spätere Auswertung. 
Die Messung wird für zwei unterschiedliche Strahlungsleistungen durchgeführt. Die eingestrahlte Leistung wird aus der maximalen Lampenleistung von 120 Wund der eingestellten Spannung am Transformator abgeschätzt, dessen maximale Spannung 250 V beträgt. Im Versuch wurden Messreihen bei 40\% und 80\% der maximalen Spannung aufgenommen, um den Einfluss der Beleuchtungsstärke auf das Verhalten der Solarzelle zu untersuchen. Um Messfehler zu minimieren, werden sämtliche Messungen unter möglichst vollständigem Ausschluss von Fremdlicht durchgeführt und der Dunkelstrom aufgenommen, um systematische Einflüsse berücksichtigen zu können. 

\begin{figure}[H]
    \centering
    
    \begin{minipage}[t]{0.45\linewidth}
        \centering
        \includegraphics[width=\linewidth]{Hydrocar1.jpeg}
        \caption{Versuchsaufbau der Elektrolyse des Hydro-Cars mittels Solarstrom.}
    \end{minipage}
    \hspace{0.05\linewidth}
    \begin{minipage}[t]{0.45\linewidth}
        \centering
        \includegraphics[width=\linewidth]{Hydrocar2.jpeg}
        \caption{Hydro-Car mit angeschlossener Brennstoffzelle.}
    \end{minipage}
    
\end{figure}
\newpage
